% ============================================================================
% II. Причини за загуба и повреда
% ============================================================================
\section{Причини за загуба и повреда на видеофайлове}

Методът за възстановяване се избира според това какво точно се е случило с файла.
Различните причини оставят различни следи върху носителя. По-долу те са групирани в
пет случая.

\subsection{Изтриване на файла}

При изтриване файлът обикновено не се унищожава. Файловата система само отбелязва
като изтрит записа, който сочи къде се намира файлът. Данните остават на носителя, но
операционната система вече смята мястото за свободно и може да го презапише по всяко
време \cite{nist80086}. На магнитен твърд диск данните остават физически на носителя, докато
мястото не бъде използвано за нов запис. Доколко могат да се възстановят и метаданните
(име, размер, списък на клъстерите), зависи от конкретната файлова система и от това колко
време е работил носителят след изтриването.

При твърдотелните дискове (SSD) това предположение не е сигурно. Bell и Boddington
показват експериментално, че вграденият софтуер на SSD може сам да изтрие
освободените блокове чрез т.нар. \term{събиране на отпадъци} (garbage collection), без
команда от компютъра \cite{bell2010}. Тогава изтритият видеозапис може да е унищожен
още преди носителят да бъде иззет.

\subsection{Форматиране на носителя}

Бързото форматиране създава нови, празни структури на файловата система. На магнитен
твърд диск то не изтрива самите данни: в опита на Bell и Boddington бързо форматиран
твърд диск запазва всички файлове \cite{bell2010}. Връзката между файловете и техните клъстери обаче
се губи. Тогава остава единствено търсенето на файлове по тяхното съдържание, без
помощта на файловата система (раздел~3.2). Ограничението от раздел~2.1 важи и тук. В същия опит
бързо форматиран SSD изчиства почти всички данни сам, за около три минути
\cite{bell2010}.

\subsection{Прекъснат запис}

Този случай е специфичен за видеото. Устройството може да пише файла MP4
последователно: първо \code{ftyp}, после растящата кутия \code{mdat}, а кутията
\code{moov} едва при нормално спиране на записа, когато таблиците на пробите са вече
известни. Ако при такава подредба записът бъде прекъснат (изтощена батерия, отнето
захранване, извадена карта, повреда на устройството), \code{moov} не се записва. Резултатът е файл с пълни видеоданни, но без метаданни.
Обикновените плейъри не могат да го отворят, въпреки че почти цялото съдържание е
налице. За възстановяване на такива прекъснати записи е създаден инструментът \code{untrunc}
\cite{untrunc},
разгледан в раздел~3.5. Фрагментираният вариант на MP4 (раздел~1.1) намалява загубата,
защото всеки завършен фрагмент носи собствените си метаданни.

\subsection{Фрагментиране на файла върху носителя}

Файловата система не винаги разполага файла в последователни клъстери. Когато свободното
място е разпокъсано, един файл може да бъде разделен на две или повече несъседни части
\cite{garfinkel2007}. Докато метаданните на файловата система са налице, това не е проблем.
След изтриване или форматиране обаче фрагментираният файл не може да бъде възстановен
чрез просто копиране на байтовете от началото до края. Garfinkel изследва над 300 употребявани
носителя и намира, че около 6\,\% от файловете на тях са фрагментирани. Разпределението
е силно неравномерно: на около половината носители няма нито един фрагментиран файл, а
на 30 носителя делът надхвърля 10\,\% \cite{garfinkel2007}. Самият автор определя
фрагментирането като сравнително рядко, но смята възстановяването на фрагментирани файлове
за необходима възможност на всеки инструмент \cite{garfinkel2007}.

\subsection{Частично презаписване}

Когато част от клъстерите на изтрит файл е използвана за нови данни, съответната част
от записа е загубена безвъзвратно. Никой метод не може да възстанови презаписаните
байтове. Възможно е обаче да се спаси останалото. Na и съавтори показват, че при
повреден видеофайл с презаписани участъци по-голямата част от съдържанието в
непрезаписаните участъци може да бъде възстановена \cite{na2014}. Тук значение има
структурата на H.264 от раздел~1.2: оцелелите кадри могат да се декодират, ако е оцелял
и кадърът, от който зависят.

Таблица~\ref{tab:causes} обобщава какво остава на носителя във всеки от петте случая.

\begin{table}[H]
\caption{Какво остава на носителя при различните причини за загуба}
\label{tab:causes}
\centering
\begin{tabular}{@{}L{3.6cm}L{3.4cm}L{3.0cm}L{3.6cm}@{}}
\toprule
\textbf{Причина} & \textbf{Метаданни на файловата система} & \textbf{Кутия \code{moov}} & \textbf{Видеоданни} \\
\midrule
Изтриване & Частично & Налична & Налични, ако не са презаписани \\
Форматиране & Загубени & Налична & Налични, ако не са презаписани \\
Прекъснат запис & Налични & Липсва & Налични \\
Фрагментиране & Загубени (след изтриване) & Може да е в друг фрагмент & Разпокъсани \\
Презаписване & Без значение & Може да е загубена & Частично загубени \\
\bottomrule
\end{tabular}
\end{table}
